\textbf{Double descent.} Belkin et al.~\cite{belkin2018reconciling} and Nakkiran et al.~\cite{nakkiran2019deep} are the empirical references; the latter varies width, epochs and label noise in deep networks. Our work uses a model small enough to sweep exactly through the threshold. \textbf{Random features theory.} Mei and Montanari~\cite{mei2019generalization} give the precise asymptotics of random-features ridge regression, d'Ascoli et al.~\cite{dascoli2020double} the bias--variance decomposition in the lazy regime, Hastie et al.~\cite{hastie2019surprises} ridgeless interpolation, Meng et al.~\cite{meng2022multiple} multiple descent when feature types are concatenated, and Liu et al.~\cite{liu2021double} random features trained with SGD. These are asymptotic or non-asymptotic theoretical results for idealised data; we measure finite $n=200$ with a nonlinear teacher and a real dataset. \textbf{Regularisation and the peak.} Tsigler and Bartlett~\cite{tsigler2020benign} study ridge regression in the over-parameterised linear setting where interpolating noisy data can still generalise. Nakkiran et al.~\cite{nakkiran2020optimal} study whether optimal regularisation avoids double descent, proving monotonicity for certain isotropic linear models. Yilmaz and Heckel~\cite{yilmaz2022regularization} show that the risk can itself be double-descent shaped as a function of the regularisation strength, which is relevant to our ridge sweep. Kan et al.~\cite{kan2020avoiding} avoid the peak in random-feature models with a hybrid of early stopping and weight decay. What we add is a side-by-side comparison, with a single hump statistic, of fixed, globally tuned, per-width validation-tuned and leave-one-out-tuned ridge on the same random-features sweep, including a failure analysis.
